\chapter{Introduction}

\section{Problem Description and Motivation}
\label{sec:problem}

The exponential growth of emerging technologies---artificial intelligence, machine learning, distributed systems, and advanced computing---has produced technical knowledge fragmented across heterogeneous platforms: scholarly repositories (ArXiv), community forums (Hacker News), developer blogs (DEV.to), and open-source ecosystems (GitHub). This fragmentation makes timely synthesis increasingly difficult for technology professionals.

In practical terms, a developer tracking the latest LLM frameworks or retrieval strategies may need to manually monitor multiple platforms, reconcile conflicting claims, and verify whether recommendations are grounded in credible sources. This manual workflow is slow, error-prone, and difficult to sustain as update velocity increases.

Purely generative large language model (LLM) assistants are insufficient for emerging technology intelligence tasks due to their static knowledge cutoff, susceptibility to hallucinated responses, lack of verifiable evidence grounding, and inability to integrate with continuously updating data streams (Ji et al., 2023). Although Retrieval-Augmented Generation (RAG) architectures partially address these limitations by incorporating external context, the relative benefit of augmenting vector-only retrieval with lexical and recency signals remains under-evaluated in cost-constrained, single-domain settings. Moreover, baseline vector-only systems typically lack temporal relevance signals, structured metadata constraints, and recency-aware ranking mechanisms, which are important for accurately capturing rapidly evolving technology trends.

A further gap exists at the systems level: many RAG research prototypes focus primarily on retrieval accuracy without addressing production-readiness---deterministic pipeline design, idempotent ingestion, monitoring and observability, deployment automation, and cloud cost governance are frequently neglected, limiting their applicability in real-world operational environments.

While prior RAG studies explore hybrid retrieval variants, most do not jointly evaluate retrieval quality, latency, cost, and production deployment constraints within one unified framework. This fragmentation in evaluation and system scope limits practical adoption, especially for small teams operating under strict budget and operational constraints. This motivates the need for an automated, continuously updating pipeline that ensures scalability, consistency, and reproducibility in technology intelligence extraction.

The central engineering and research challenge addressed in this project is therefore:

\begin{quote}
\textbf{\textit{How can a scalable, continuously updating hybrid retrieval system be designed and deployed to synthesize emerging technology trends from heterogeneous scholarly and industry sources while satisfying latency and cost constraints?}}
\end{quote}

This project is motivated by three converging factors. First, technology professionals increasingly require timely and citation-grounded intelligence that synthesizes insights from both scholarly research and rapidly evolving industry discussions. Second, mature serverless cloud infrastructures now enable cost-effective deployment of production-grade data and machine learning pipelines. Third, a gap remains in the Retrieval-Augmented Generation (RAG) literature between retrieval-focused research prototypes and operationally robust systems capable of continuous ingestion, evaluation, and deployment in real-world environments.

\section{Research Questions}
\label{sec:rq}

To address the central problem, this study investigates the following research questions:

\begin{enumerate}[leftmargin=0.5in]
  \item \textbf{Data Engineering.} How can a cloud-native pipeline reliably ingest heterogeneous scholarly and industry data in a deterministic, idempotent manner?

  \item \textbf{Retrieval Optimization.} How does hybrid retrieval (metadata filtering, BM25, and Reciprocal Rank Fusion) compare to baseline vector-only retrieval in context precision and faithfulness?

  \item \textbf{Operational Trade-offs.} What trade-offs arise between retrieval precision, p95 latency, and cost per query?

  \item \textbf{MLOps Integration.} How can deployment automation, monitoring, and cost governance ensure the reliability of a production-grade RAG system?
\end{enumerate}

\section{Objectives of the Study}
\label{sec:objectives}

The primary objective is to design, implement, and evaluate a hybrid Retrieval-Augmented Generation system capable of synthesizing real-time emerging technology intelligence within a scalable, cost-aware cloud-native architecture. In this project, machine learning model development is approached from a systems perspective: rather than training a large model from scratch, the work focuses on selecting appropriate pretrained models, integrating them into a retrieval-augmented architecture, tuning retrieval and ranking behavior through systematic ablation, and evaluating quality, latency, and cost trade-offs in a production-like deployment setting. The specific objectives are:

\begin{enumerate}[leftmargin=0.5in]
  \item To construct a cloud-native data engineering pipeline for batch and near-real-time ingestion from five heterogeneous data sources.
  \item To implement and compare baseline vector-only retrieval against hybrid retrieval with metadata filtering, full-corpus BM25, and Reciprocal Rank Fusion (RRF).
  \item To evaluate system performance across context precision, faithfulness, p95 retrieval latency, and cost per query.
  \item To deploy and monitor the system using serverless AWS infrastructure consistent with MLOps best practices (Sculley et al., 2015; Breck et al., 2017).
\end{enumerate}

\section{Significance of the Study}
\label{sec:significance}

This project contributes at the intersection of data engineering, natural language processing, and MLOps in three ways. First, it provides controlled empirical insights within a single-domain, moderate-scale evaluation setting showing that hybrid retrieval (weighted BM25+RRF with cross-encoder reranking) shows consistent improvement over vector-only baselines across faithfulness, answer relevancy, and context precision---a configuration common in practice but under-represented in the literature; these findings may not generalize to heterogeneous or large-scale multi-domain corpora without further validation; in settings where lexical diversity is higher, hybrid retrieval may yield larger gains than observed here. (Citation grounding shows a marginal decrease of $-0.02$, discussed in Chapter~\ref{ch:eval}.) Second, it demonstrates that weighted Reciprocal Rank Fusion combined with a lightweight cross-encoder (ms-marco-MiniLM-L-6-v2) achieves a composite quality gain of $+0.047$ over the baseline while maintaining a manageable latency trade-off (mean $+1.36$\,s, Cohen's $d = -0.34$, small effect). Third, and most importantly, the primary contribution of this work lies in the integration of hybrid retrieval, drift detection, hallucination verification, and cost governance operationalized together in a single deployable architecture on the AWS Free Tier, providing a practical blueprint for student and small-team RAG deployments.

\section{Alignment Between Research Questions and Evaluation}
\label{sec:rq_eval_map}

Table~\ref{tab:rq_map} provides an explicit mapping from each research question to the corresponding evaluation metric and experimental configuration used to answer it.

\begin{table}[H]
\caption{Mapping of Research Questions to Evaluation Metrics and Methodology.}
\label{tab:rq_map}
\centering
\begin{tabular}{p{0.5cm}p{4.0cm}p{3.5cm}p{4.0cm}}
\toprule
\textbf{RQ} & \textbf{Question (abbreviated)} & \textbf{Primary Metric} & \textbf{Evaluation Approach} \\
\midrule
1 & Cloud-native deterministic ingestion & Pipeline idempotency; SHA-256 deduplication pass rate & Unit tests (24 ingestion tests); repeated pipeline runs \\[4pt]
2 & Hybrid vs.\ baseline retrieval precision & Context precision (RAGAS); faithfulness; answer relevancy; composite score; citation grounding $> 0.80$ & Paired comparison over 50 queries; Wilcoxon signed-rank test; Cohen's $d$; bootstrap 95\% CI \\[4pt]
3 & Latency and cost trade-offs & p95 latency $< 2$s (retrieval stage); cost $<$ USD~15/month & Per-query timing; token usage profiling; cost projection \\[4pt]
4 & MLOps integration & Drift detection coverage; CI/CD pass rate; health-check uptime & Drift simulation runs; CloudWatch alarm verification \\
\bottomrule
\end{tabular}
\medskip
\noindent\small\textit{Note: The citation grounding threshold of 0.80 represents the evaluation target under RAGAS scoring, while a lower operational threshold (0.25) is used at runtime to ensure robustness under constrained LLM performance.}
\end{table}

\medskip
\noindent The study is structured as follows: Chapter~2 reviews the relevant literature; Chapter~3 presents the methodology and hypotheses; Chapter~4 describes the system architecture; Chapter~5 defines the evaluation framework and presents results; Chapter~6 documents implementation details; and Chapter~7 presents conclusions, limitations, and future work. The following chapter reviews the technical foundations underpinning these design decisions.
